\section{DECLARAÇÃO SOBRE USO DE INTELIGÊNCIA ARTIFICIAL }

Em conformidade com as diretrizes de integridade acadêmica e transparência na pesquisa científica, o autor declara o uso de ferramentas de Inteligência Artificial Generativa — especificamente o modelo de linguagem Google Gemini — como tecnologia assistiva durante as etapas de ideação, redação, formatação e revisão deste Trabalho de Conclusão de Curso.

O emprego da referida ferramenta ocorreu sob supervisão e direcionamento estrito do autor, restringindo-se às seguintes frentes operacionais:
\begin{itemize}[topsep=0pt, noitemsep]
    \item \textbf{Geração Direta de Texto:} A ferramenta foi empregada para a consolidação textual do Resumo (e sua respectiva tradução para o \textit{Abstract}) e para a redação desta própria seção de Declaração de Uso de Inteligência Artificial, bem como a redação dos textos introdutórios, termos de consentimento e cabeçalhos de instrução do formulário de pesquisa. Destaca-se que a geração ocorreu estritamente a partir das premissas, dados analíticos e direcionamentos estruturais fornecidos previamente pelo autor.

    \item \textbf{Consultoria Metodológica, Normativa e Ética:} Realização de consultas referentes à conformidade do documento com as regras da Associação Brasileira de Normas Técnicas (ABNT), abrangendo o refatoramento de citações e a obtenção de guias orientativos sobre o escopo de conteúdo exigido para determinadas seções do trabalho. O auxílio estendeu-se à redação estratégica para a anonimização de dados organizacionais sensíveis, assegurando o sigilo e a confidencialidade ao longo do relato de experiência.

    \item \textbf{Revisão Gramatical e Refinamento Estilístico:} Lapidação do texto em língua portuguesa, auxiliando na adequação da sintaxe, eliminação de vícios de linguagem, garantia de coesão textual e elevação do vocabulário para manter o tom formal, objetivo e impessoal exigido pela escrita acadêmica.

    \item \textbf{Codificação LaTeX e Visualização de Dados:} A ferramenta atuou como suporte técnico na formatação estrutural do documento e na resolução de problemas de compilação, assegurando a conformidade com as normas da ABNT. O auxílio abrangeu a organização modular do código-fonte, o encapsulamento de arquivos PDF externos na seção de apêndices e a elaboração de \textit{scripts} de visualização de dados utilizando pacotes específicos (como \textit{pgfplots} e \textit{pgf-pie}). Esse suporte permitiu a renderização de representações geométricas adequadas para dados categóricos --- a exemplo de gráficos de barras horizontais empilhadas para escalas \textit{Likert} --- e a definição de paletas de cores e legendas alinhadas semanticamente aos resultados da pesquisa.

    \item \textbf{Discussão Conceitual e Ideação:} Uso do modelo de linguagem como ferramenta de \textit{brainstorming}, englobando a formulação de ideias de títulos para o trabalho e a validação de conceitos teóricos durante a fase de modelagem --- a exemplo do debate analítico sobre a definição de Casos de Uso com base na entrega de valor de negócio em contraposição a fluxos estritamente navegacionais e o debate sobre a usabilidade do instrumento de coleta de dados, definindo a melhor estratégia de apresentação da documentação para os participantes da pesquisa.

    \item \textbf{Estruturação do Instrumento de Pesquisa:} O modelo auxiliou na elaboração do formulário de coleta de dados aplicado à equipe de desenvolvimento. O suporte englobou a estruturação de perguntas para medir a percepção de implementação frente à documentação de requisitos, a definição das escalas de avaliação e a redação do termo de consentimento. Adicionalmente, a ferramenta auxiliou na refatoração do documento de especificação anexo (CdUs, RNs e Dicionários de Dados), organizando-o em um formato de consulta otimizado para reduzir a carga cognitiva dos respondentes durante a pesquisa.

    \item \textbf{Apoio na Comunicação com Participantes da Pesquisa:} A ferramenta foi empregada para redigir, revisar e refinar as mensagens de convite, acompanhamento e encerramento enviadas à equipe de desenvolvimento e coordenação do SFM. O uso consistiu em adequar o tom e a clareza da comunicação para engajar os profissionais a responderem ao instrumento de coleta de dados (questionário eletrônico), otimizando a taxa de resposta da pesquisa.
\end{itemize}
    
Por fim, o autor reitera que a utilização destas ferramentas de Inteligência Artificial possuiu caráter consultivo e de apoio estrutural. Dessa forma, o autor atesta ter revisado integralmente o material gerado e assume a responsabilidade total, irrestrita e exclusiva sobre todo o conteúdo intelectual, analítico e técnico apresentado neste documento, bem como pela originalidade da pesquisa e pela garantia de ausência de plágio.